The number $993$ is fortunate, since its digit sum is $21$. The numbers $994,\ldots,999$ have digit sums $22,\ldots,27$, and $1000,\ldots,1005$ have digit sums $1,\ldots,6$. None is divisible by $7$, so $993$ is superfortunate.

To prove minimality, let a fortunate number $n<993$ end in digit $r$. If $r\le2$, adding $7$ causes no carry and gives another fortunate number. If $r\ge4$, among the next twelve numbers there are at least seven consecutive numbers after the first tens boundary and before the next. Their digit sums are seven consecutive integers, so one is fortunate.

It remains to consider $r=3$. The first tens boundary occurs at $n+7$. Let $c$ be the number of carried positions there. Since $n<993$ and ends in $3$, $c$ is either $1$ or $2$. For $7\le j\le12$, the digit sum changes by $j-9c$. Modulo $7$ this is $j-2c$. If $c=1$, choose $j=9$; if $c=2$, choose $j=11$. In either case another fortunate number occurs within twelve steps. Thus no smaller positive integer works, and the answer is $\boxed{993}$.
